% ============================================================
%  C: Como Programar -- 6ª Edição | Deitel & Deitel
%  Capítulo 14 -- Outros Tópicos sobre C
%  EXERCÍCIO DE AUTORREVISÃO (14.1)
%  Abordagem incremental: Caps. 1 a 14
% ============================================================
\documentclass[12pt,a4paper]{article}
\usepackage[utf8]{inputenc}
\usepackage[T1]{fontenc}
\usepackage[brazil]{babel}
\usepackage{geometry}
\geometry{top=2.5cm,bottom=2.5cm,left=3cm,right=3cm}
\usepackage{parskip}\setlength{\parskip}{6pt}\setlength{\parindent}{0pt}
\usepackage{xcolor,tcolorbox,enumitem,listings,fancyhdr,titlesec,hyperref,booktabs,array}
\tcbuselibrary{skins,breakable}

\definecolor{deitelblue}{RGB}{0,70,127}
\definecolor{deitelgray}{RGB}{240,242,245}
\definecolor{deitelgreen}{RGB}{0,110,60}
\definecolor{deitellightblue}{RGB}{220,235,250}
\definecolor{deitelred}{RGB}{160,20,20}
\definecolor{deitellorange}{RGB}{200,90,0}

\newtcolorbox{boaprática}[1][]{colback=deitellightblue,colframe=deitelblue,
  fonttitle=\bfseries\small,title={Boa prática de programação},breakable,#1}
\newtcolorbox{erroco}[1][]{colback=red!5,colframe=deitelred,
  fonttitle=\bfseries\small,title={Erro comum de programação},breakable,#1}
\newtcolorbox{respostabox}[1][]{colback=deitelgray,colframe=deitelblue!60,
  fonttitle=\bfseries\small\color{deitelblue},title={Resposta detalhada},breakable,#1}
\newtcolorbox{enunciadoBox}[1][]{colback=white,colframe=deitelblue,
  fonttitle=\bfseries,title={#1},breakable}
\newtcolorbox{saida}[1][]{colback=black!88,colframe=deitelblue,
  fontupper=\ttfamily\small\color{green!70},title={\small\color{white}Saída do programa},breakable,#1}

\setlist[enumerate,1]{label=\textbf{\alph*)},leftmargin=2em}
\setlist[itemize]{leftmargin=2em}

\lstset{
  basicstyle=\ttfamily\small,backgroundcolor=\color{deitelgray},
  frame=single,framerule=0.4pt,rulecolor=\color{deitelblue!40},
  breaklines=true,showstringspaces=false,
  keywordstyle=\color{deitelblue}\bfseries,
  commentstyle=\color{deitelgreen}\itshape,
  stringstyle=\color{deitelred},
  language=C,numbers=left,numberstyle=\tiny\color{gray},
  xleftmargin=18pt,xrightmargin=5pt,stepnumber=1,
}

\pagestyle{fancy}\fancyhf{}
\fancyhead[L]{\small\color{deitelblue}\textbf{C: Como Programar -- 6ª ed. | Deitel \& Deitel}}
\fancyhead[R]{\small\color{deitelblue}Cap. 14 -- Autorrevisão}
\fancyfoot[C]{\small\thepage}
\renewcommand{\headrulewidth}{0.4pt}

\titleformat{\section}[block]{\large\bfseries\color{deitelblue}}{\thesection.}{0.5em}{}[\titlerule]
\titleformat{\subsection}[block]{\normalsize\bfseries\color{deitelblue!80}}{}{}{}

\hypersetup{colorlinks=true,linkcolor=deitelblue}

\begin{document}

\begin{titlepage}
  \centering\vspace*{2cm}
  {\color{deitelblue}\rule{\linewidth}{2pt}}\\[0.4cm]
  {\huge\bfseries\color{deitelblue}C: Como Programar}\\[0.2cm]
  {\large\color{deitelblue}6ª Edição $\cdot$ Paul Deitel \& Harvey Deitel}\\[0.2cm]
  {\color{deitelblue}\rule{\linewidth}{2pt}}\\[1cm]
  {\LARGE\bfseries Capítulo 14}\\[0.3cm]
  {\Large\color{deitelblue!80}Outros Tópicos sobre C}\\[0.8cm]
  \begin{tcolorbox}[colback=deitellightblue,colframe=deitelblue,width=0.85\linewidth]
    \centering{\large\bfseries Exercício de Autorrevisão (14.1)}\\[0.2cm]
    {\normalsize 19 lacunas conceituais\\
    Abordagem incremental: Capítulos 1 a 14\\[0.2cm]
    \textbf{O capítulo final do núcleo de C}}
  \end{tcolorbox}
\end{titlepage}

\section*{Construindo sobre os Capítulos Anteriores}

O Capítulo~14 -- ``Outros Tópicos sobre C'' -- amarra os últimos fios da
linguagem, completando seu conjunto de ferramentas para programas profissionais
em C. Construindo sobre os Caps.~1--13, agora incorporamos:

\begin{itemize}[noitemsep]
  \item \textbf{Redirecionamento de E/S} -- \texttt{<}, \texttt{>}, \texttt{>>},
  \texttt{|} (pipe) do shell
  \item \textbf{Argumentos variádicos} -- funções como \texttt{printf} que aceitam
  número variável de argumentos via \texttt{<stdarg.h>}: \texttt{va\_list},
  \texttt{va\_start}, \texttt{va\_arg}, \texttt{va\_end}
  \item \textbf{Argumentos de linha de comando} -- \texttt{int main(int argc, char *argv[])}
  \item \textbf{Compilação multi-arquivo} -- \texttt{extern}, vinculação, ferramenta
  \texttt{make} e \texttt{Makefile}
  \item \textbf{Término de programa} -- \texttt{exit}, \texttt{atexit}, \texttt{abort}
  \item \textbf{Sufixos numéricos} -- \texttt{L}, \texttt{U}, \texttt{F} para
  controlar tipo de constantes
  \item \textbf{Tratamento de sinais} -- \texttt{<signal.h>}, \texttt{signal},
  \texttt{raise}, \texttt{SIGINT}, \texttt{SIGABRT}, \ldots
  \item \textbf{Alocação dinâmica avançada} -- \texttt{calloc} (com zeros),
  \texttt{realloc} (redimensiona)
  \item \textbf{Arquivos temporários} -- \texttt{tmpfile}
  \item \textbf{Qualificadores adicionais} -- \texttt{const}, \texttt{volatile}
  \item \textbf{\texttt{goto}} -- desvio incondicional (uso restrito)
\end{itemize}

\begin{boaprática}
  Este é o capítulo que separa programadores casuais de programadores
  profissionais. Argumentos de linha de comando habilitam ferramentas Unix-style;
  \texttt{Makefile} viabiliza projetos grandes; \texttt{calloc}/\texttt{realloc}
  são essenciais para estruturas dinâmicas robustas; \texttt{signal} permite
  programas resilientes.
\end{boaprática}

\tableofcontents\newpage

% ============================================================
\section{Exercício 14.1 -- Preenchimento de Lacunas (19 itens)}
% ============================================================

\begin{respostabox}
\begin{enumerate}[label=\textbf{\alph*)}]

\item Símbolo que redireciona entrada de arquivo (em vez do teclado):
\textbf{\texttt{<}} (redirecionamento de entrada).
\begin{lstlisting}
./meuprog < entrada.txt
\end{lstlisting}

\item Símbolo que redireciona saída para arquivo (em vez da tela):
\textbf{\texttt{>}} (redirecionamento de saída).
\begin{lstlisting}
./meuprog > saida.txt
\end{lstlisting}

\item Símbolo que acrescenta saída ao final de um arquivo:
\textbf{\texttt{>>}} (acréscimo de saída).
\begin{lstlisting}
./meuprog >> log.txt
\end{lstlisting}
Diferença crucial: \texttt{>} \textit{trunca} o arquivo; \texttt{>>}
\textit{preserva} o conteúdo existente e acrescenta.

\item Direciona saída de um programa para entrada de outro:
\textbf{\texttt{|}} (pipe).
\begin{lstlisting}
ls | sort | uniq
\end{lstlisting}
A saída de \texttt{ls} torna-se entrada de \texttt{sort}, cuja saída
torna-se entrada de \texttt{uniq}. Essa composição é a essência da
filosofia Unix.

\item Lista de parâmetros indicando número variável de argumentos:
\textbf{\texttt{...}} (reticências).
\begin{lstlisting}
int media( int contagem, ... );
\end{lstlisting}
\texttt{printf} e \texttt{scanf} são variádicas.

\item Macro chamada \textit{antes} de acessar argumentos variáveis:
\textbf{\texttt{va\_start}}.
\begin{lstlisting}
va_list args;
va_start( args, contagem );  /* contagem = ultimo arg nomeado */
\end{lstlisting}

\item Macro que acessa argumentos individuais da lista variável:
\textbf{\texttt{va\_arg}}.
\begin{lstlisting}
int valor = va_arg( args, int );  /* le proximo argumento */
\end{lstlisting}

\item Macro que finaliza o processamento da lista variável:
\textbf{\texttt{va\_end}}.
\begin{lstlisting}
va_end( args );  /* SEMPRE chamar antes do return! */
\end{lstlisting}

\item Argumento de \texttt{main} que recebe \textit{quantidade} de argumentos
da linha de comando: \textbf{\texttt{argc}} (\textit{argument count}).

\item Argumento de \texttt{main} que armazena os argumentos como strings:
\textbf{\texttt{argv}} (\textit{argument vector}).
\begin{lstlisting}
int main( int argc, char *argv[] )
{
    /* argv[0] = nome do programa
       argv[1], argv[2], ... = argumentos */
}
\end{lstlisting}

\item Utilitário do Linux/Unix que processa arquivo de regras:
\textbf{\texttt{make}}, que lê um \textbf{\texttt{Makefile}}.
\begin{lstlisting}
# Exemplo de Makefile
programa: main.o util.o
    gcc -o programa main.o util.o

main.o: main.c util.h
    gcc -c main.c
\end{lstlisting}

\item Função que força término da execução: \textbf{\texttt{exit}}.
\begin{lstlisting}
exit( 0 );           /* terminacao bem-sucedida */
exit( EXIT_FAILURE ); /* terminacao com erro */
\end{lstlisting}

\item Função que registra rotina a ser chamada no término normal:
\textbf{\texttt{atexit}}.
\begin{lstlisting}
void cleanup( void ) { printf( "Limpando...\n" ); }
int main( void ) {
    atexit( cleanup );  /* chamada automatica ao final */
    return 0;
}
\end{lstlisting}

\item Letra anexada a constantes para especificar tipo exato:
\textbf{sufixo}.
\begin{lstlisting}
3.14F      /* float (sem F seria double) */
100UL      /* unsigned long */
9999999L   /* long int */
1.5L       /* long double */
\end{lstlisting}

\item Função que abre arquivo temporário (existe até fechar ou fim do programa):
\textbf{\texttt{tmpfile}}.
\begin{lstlisting}
FILE *fp = tmpfile();  /* arquivo binario, automaticamente removido */
\end{lstlisting}

\item Função para interceptar eventos inesperados: \textbf{\texttt{signal}}.
\begin{lstlisting}
void tratador( int sig ) { /* trata o sinal */ }
signal( SIGINT, tratador );  /* registra handler para Ctrl+C */
\end{lstlisting}

\item Função que gera sinal de dentro do programa: \textbf{\texttt{raise}}.
\begin{lstlisting}
raise( SIGABRT );  /* gera SIGABRT explicitamente */
\end{lstlisting}

\item Função que aloca array e \textit{inicializa em zero}:
\textbf{\texttt{calloc}}.
\begin{lstlisting}
int *arr = calloc( 100, sizeof( int ) );
/* 100 inteiros, todos zero */
\end{lstlisting}
Diferença de \texttt{malloc}: \texttt{calloc} aceita 2 argumentos (n e
tamanho) e \textit{garante zeros}; \texttt{malloc} deixa memória com lixo.

\item Função que muda tamanho de bloco previamente alocado:
\textbf{\texttt{realloc}}.
\begin{lstlisting}
arr = realloc( arr, 200 * sizeof( int ) );
/* expande para 200 inteiros, preserva os 100 antigos */
\end{lstlisting}

\end{enumerate}
\end{respostabox}

\newpage

% ============================================================
\section*{Resumo Visual das Funções Novas do Cap. 14}

\begin{tcolorbox}[colback=deitelgray,colframe=deitelblue,title={\bfseries Mapa de funções por categoria}]

\textbf{Argumentos variádicos (\texttt{<stdarg.h>}):}
\begin{center}
\renewcommand{\arraystretch}{1.2}
\begin{tabular}{l l}
\toprule
\textbf{Macro} & \textbf{Função} \\ \midrule
\texttt{va\_list}  & tipo para apontar para argumentos \\
\texttt{va\_start} & inicializa va\_list \\
\texttt{va\_arg}   & obtém próximo argumento \\
\texttt{va\_end}   & finaliza (sempre chamar!) \\
\bottomrule
\end{tabular}
\end{center}

\textbf{Término de programa (\texttt{<stdlib.h>}):}
\begin{center}
\begin{tabular}{l l}
\toprule
\texttt{exit}    & encerra programa, retorna código \\
\texttt{atexit}  & registra função para chamar no final \\
\texttt{abort}   & término anormal (gera SIGABRT) \\
\bottomrule
\end{tabular}
\end{center}

\textbf{Alocação dinâmica (\texttt{<stdlib.h>}):}
\begin{center}
\begin{tabular}{l l}
\toprule
\texttt{malloc}  & aloca bloco, conteúdo indefinido \\
\texttt{calloc}  & aloca array, inicializa zero \\
\texttt{realloc} & redimensiona bloco \\
\texttt{free}    & libera bloco \\
\bottomrule
\end{tabular}
\end{center}

\textbf{Sinais (\texttt{<signal.h>}):}
\begin{center}
\begin{tabular}{l l}
\toprule
\texttt{signal} & registra tratador para sinal \\
\texttt{raise}  & gera sinal de dentro do programa \\
\bottomrule
\end{tabular}
\end{center}

\textbf{Arquivos:}
\begin{center}
\begin{tabular}{l l}
\toprule
\texttt{tmpfile}  & cria arquivo temporário \\
\texttt{tmpnam}   & gera nome único para arquivo temporário \\
\bottomrule
\end{tabular}
\end{center}

\end{tcolorbox}

% ============================================================
\section*{Gabarito Rápido}

\begin{tcolorbox}[colback=deitellightblue,colframe=deitelblue,
  title={\bfseries\large Respostas Resumidas -- 14.1}]
\begin{itemize}[noitemsep]
  \item[a)] \texttt{<} \quad b) \texttt{>} \quad c) \texttt{>>} \quad d) pipe \texttt{|}
  \item[e)] reticências \texttt{...} \quad f) \texttt{va\_start} \quad g) \texttt{va\_arg}
  \item[h)] \texttt{va\_end} \quad i) \texttt{argc} \quad j) \texttt{argv}
  \item[k)] \texttt{make}, \texttt{Makefile} \quad l) \texttt{exit} \quad m) \texttt{atexit}
  \item[n)] sufixo \quad o) \texttt{tmpfile} \quad p) \texttt{signal}
  \item[q)] \texttt{raise} \quad r) \texttt{calloc} \quad s) \texttt{realloc}
\end{itemize}
\end{tcolorbox}

\end{document}
